\documentclass[../main.tex]{subfiles}

\begin{document}

\chapter{System of Linear Equations}

One of the most important application of matrices would be in
linear systems of equations. With the ideas on matrices, we can
open our third eye when viewing system of linear equations.
This note is dedicated for matrices in linear equations,
Gaussian Elimination, and LU decompositions.

\section{Basics and Intuitions}

As always, let's get started with fundamental definitions, concepts,
and intuitions that will be used throughout.

First, I am pretty sure most of you are aware of what system of
equation is. However, because it is always nice to formally state
our understanding, here is the definition of linear system of
equation.

\begin{definition}
A \textbf{system of linear equations} is composed of linear equations,
or equations whose highest degree is $1$. A system of $m$ equations
and $n$ variables adhere to the following forms where $a_{ij}$ and
$b_{i}$ are known quantities.
\begin{align*}
    a_{11}x_1 + a_{12}x_2 + a_{13}x_3 + \cdots + a_{1n}x_n &= b_1 \\
    a_{21}x_1 + a_{22}x_2 + a_{23}x_3 + \cdots + a_{2n}x_n &= b_2 \\
    a_{31}x_1 + a_{32}x_2 + a_{33}x_3 + \cdots + a_{3n}x_n &= b_3 \\
    \vdots\qquad\qquad& \\
    a_{m1}x_1 + a_{m2}x_2 + a_{m3}x_3 + \cdots + a_{mn}x_n &= b_m
\end{align*}
\end{definition}

Below is an example of a system of linear equations.
\begin{align*}
    x + 2y + 3z &= 4 \\
    -x + 3y + 2z &= 5
\end{align*}
Notice that all equations above are linear, meaning, their degrees
are $1$. Now, what does it mean to find a solution to such systems?

\begin{definition}
A \textbf{solution to a system of linear equations} is a set of
scalar values for unknown variables such that all linear equations
are satisfied when substituted.
\end{definition}

For instance, $(2, 1)$ is the solution to the following system of
linear equations.
\begin{align*}
    x - 2y &= 0 \\
    x + y  &= 3
\end{align*}
To check it, you can simply substitute to see if the equations are
satisfied.

Now, if you recall from the previous note where we discussed matrix
multiplication, system of linear equations can be written as a
matrix multiplication.

\begin{definition}
A system of linear equation
\begin{align*}
    a_{11}x_1 + a_{12}x_2 + a_{13}x_3 + \cdots + a_{1n}x_n &= b_1 \\
    a_{21}x_1 + a_{22}x_2 + a_{23}x_3 + \cdots + a_{2n}x_n &= b_2 \\
    a_{31}x_1 + a_{32}x_2 + a_{33}x_3 + \cdots + a_{3n}x_n &= b_3 \\
    \vdots\qquad\qquad& \\
    a_{m1}x_1 + a_{m2}x_2 + a_{m3}x_3 + \cdots + a_{mn}x_n &= b_m
\end{align*}
can be written as $Ax = b$ where $A$, $x$, and $b$ are matrices
defined as the following.
\[
A =
\begin{bmatrix}
    a_{11} & a_{12} & a_{13} & \cdots & a_{1n} \\
    a_{21} & a_{22} & a_{23} & \cdots & a_{2n} \\
    a_{31} & a_{32} & a_{33} & \cdots & a_{3n} \\
    \vdots & \vdots & \vdots & \ddots & \vdots \\
    a_{m1} & a_{m2} & a_{m3} & \cdots & a_{mn} \\
\end{bmatrix},
\quad
x =
\begin{bmatrix}
    x_1 \\
    x_2 \\
    x_3 \\
    \vdots \\
    x_n
\end{bmatrix},
\quad
b
\begin{bmatrix}
    b_1 \\
    b_2 \\
    b_3 \\
    \vdots \\
    b_m
\end{bmatrix}
\]
\end{definition}

When you actually perform matrix multiplication, you can see that
$Ax = b$ is identical to the system of linear equations. Now that
we have basic understanding of the relationship between matrices
and linear system of equations, let's take a look at a theorem and
an important corollary.

\begin{theorem}{Different Solutions to a System}
For two distinct solutions $x_1$ and $x_2$ of $Ax = b$,
$\alpha x_1 + \beta x_2$ is also a solution to the linear system
given that $\alpha, \beta \in \mathbb{R}$ and $\alpha + \beta = 1$.
\end{theorem}
\begin{proof}
From Theorem \ref{theo:Properties of Matrix Multiplication}, it was
shown that $\lambda (AB) = A (\lambda B)$ for scalar $\lambda$ and
matrices $A$ and $B$ that can be multiplied. Substituting $\alpha x_1
+ \beta x_2$, the following system is obtained.
\begin{align*}
    A (\alpha x_1 + \beta x_2)
    &= A(\alpha x_1) + A(\beta x_2)
    = \alpha (Ax_1) + \beta (Ax_2) \\
    &= \alpha b + \beta b
    = b
\end{align*}
Therefore, $(\alpha x_1 + \beta x_2)$ is another solution to the
system.
\end{proof}

From this theorem, we can establish an important corollary about the
number of solutions to a system of linear equations.

\begin{corollary}{Number of Solutions to a Linear System}
A system of linear equations can retain $0$, $1$, or infinitely many
solutions.
\end{corollary}
\begin{proof}
First, it is self-evident that a system can retain $0$ or $1$
solution. Therefore, the fact that there exists infinite solutions
if there are more than $1$ solution must be proven.

Let $f(\alpha) = \alpha x_1 + (1 - \alpha) x_2$ be a function that
returns possible solutions where $x_1 \neq x_2$ are fixed values.
It suffices to show that the function $f$ is injective as there exists
infinitely many values of $\alpha$. Consider the following equations.
\begin{align*}
    f(\alpha)
    &= \alpha x_1 + (1 - \alpha) x_2 
     = \alpha (x_1 - x_2) + x_2 \\
    f(\alpha')
    &= \alpha' x_1 + (1 - \alpha') x_2 
     = \alpha' (x_1 - x_2) + x_2
\end{align*}
If $f(\alpha) = f(\alpha')$, then $\alpha (x_1 - x_2) + x_2 =
\alpha' (x_1 - x_2) + x_2$ and $\alpha = \alpha'$, proving the
injectivity. Therefore, a system of linear equations can only have
$0$, $1$, or infinitely many solutions.
\end{proof}

In middle school, you probably saw a visual understanding of the
corollary above. Let's do the same, except in 3-dimensional space.
Below are two examples of when there exist no solution. Note that
not all three planes intersect at a single point.

\begin{minipage}{0.48\textwidth}
\tdplotsetmaincoords{80}{100}
\begin{tikzpicture}[
        tdplot_main_coords,
        axis/.style={->,thick},
        scale=0.40
    ]
    \pgfmathsetmacro{\a}{1.13397459622}
    \pgfmathsetmacro{\b}{6.33012701892}

    \foreach \x in {0,1,...,9}
        \foreach \y in {0,1,...,9} {
			\draw[gray,very thin]
                (\x,0,0) -- (\x,9.5,0);
			\draw[gray,very thin] (0,\y,0) -- (9.5,\y,0);
		}
    \foreach \y in {0,1,...,9}
        \foreach \z in {0,1,...,9} {
			\draw[gray,very thin]
                (0,\y,0) -- (0,\y,9.5);
			\draw[gray,very thin] (0,0,\z) -- (0,9.5,\z);
		}
    \foreach \z in {0,1,...,9}
        \foreach \x in {0,1,...,9} {
			\draw[gray,very thin]
                (0,0,\z) -- (9.5,0,\z);
			\draw[gray,very thin] (\x,0,0) -- (\x,0,9.5);
		}

	\draw[axis] (-2,0,0) -- (10,0,0) node[left] {$x$};
	\draw[axis] (0,-1,0) -- (0,10,0) node[right] {$y$};
	\draw[axis] (0,0,-1) -- (0,0,10) node[left] {$z$};

	\draw[opacity=.5,very thick,fill=Red]
        (2,2,2) -- (8,2,2) -- (8,8,2) -- (2,8,2) -- cycle;
	\draw[opacity=.5,very thick,fill=Blue]
        (2,2.5,\a) -- (8,2.5,\a) -- (8,5.5,\b) -- (2,5.5,\b) -- cycle;
	\draw[opacity=.5,very thick,fill=yellow]
        (2,7.5,\a) -- (8,7.5,\a) -- (8,4.5,\b) -- (2,4.5,\b) -- cycle;
    \draw[very thick, dashed] (2,3,2) -- (8,3,2);
    \draw[very thick, dashed] (2,7,2) -- (8,7,2);
    \draw[very thick, dashed] (2,5,5.46410161514) --
        (8,5,5.46410161514);
\end{tikzpicture}
\end{minipage}
\hfill
\begin{minipage}{0.48\textwidth}
\tdplotsetmaincoords{80}{100}
\begin{tikzpicture}[
        tdplot_main_coords,
        axis/.style={->,thick},
        scale=0.40
    ]
    \foreach \x in {0,1,...,9}
        \foreach \y in {0,1,...,9} {
			\draw[gray,very thin]
                (\x,0,0) -- (\x,9.5,0);
			\draw[gray,very thin] (0,\y,0) -- (9.5,\y,0);
		}
    \foreach \y in {0,1,...,9}
        \foreach \z in {0,1,...,9} {
			\draw[gray,very thin]
                (0,\y,0) -- (0,\y,9.5);
			\draw[gray,very thin] (0,0,\z) -- (0,9.5,\z);
		}
    \foreach \z in {0,1,...,9}
        \foreach \x in {0,1,...,9} {
			\draw[gray,very thin]
                (0,0,\z) -- (9.5,0,\z);
			\draw[gray,very thin] (\x,0,0) -- (\x,0,9.5);
		}

	\draw[axis] (-2,0,0) -- (10,0,0) node[left] {$x$};
	\draw[axis] (0,-1,0) -- (0,10,0) node[right] {$y$};
	\draw[axis] (0,0,-1) -- (0,0,10) node[left] {$z$};

	\draw[opacity=.5,very thick,fill=Red]
        (2,2,2) -- (8,2,2) -- (8,8,2) -- (2,8,2) -- cycle;
	\draw[opacity=.5,very thick,fill=Blue]
        (2,2,4) -- (8,2,4) -- (8,8,4) -- (2,8,4) -- cycle;
	\draw[opacity=.5,very thick,fill=Yellow]
        (2,2,6) -- (8,2,6) -- (8,8,6) -- (2,8,6) -- cycle;
\end{tikzpicture}
\end{minipage}

Continuing, the left diagram represents the case where there
is a single solution while the right diagram shows the system
with infinite solutions.

\begin{minipage}{0.48\textwidth}
\tdplotsetmaincoords{80}{100}
\begin{tikzpicture}[
        tdplot_main_coords,
        axis/.style={->,thick},
        scale=0.40
    ]
    \pgfmathsetmacro{\a}{6.24264068712}

    \foreach \x in {0,1,...,9}
        \foreach \y in {0,1,...,9} {
			\draw[gray,very thin]
                (\x,0,0) -- (\x,9.5,0);
			\draw[gray,very thin] (0,\y,0) -- (9.5,\y,0);
		}
    \foreach \y in {0,1,...,9}
        \foreach \z in {0,1,...,9} {
			\draw[gray,very thin]
                (0,\y,0) -- (0,\y,9.5);
			\draw[gray,very thin] (0,0,\z) -- (0,9.5,\z);
		}
    \foreach \z in {0,1,...,9}
        \foreach \x in {0,1,...,9} {
			\draw[gray,very thin]
                (0,0,\z) -- (9.5,0,\z);
			\draw[gray,very thin] (\x,0,0) -- (\x,0,9.5);
		}

	\draw[axis] (-2,0,0) -- (10,0,0) node[left] {$x$};
	\draw[axis] (0,-1,0) -- (0,10,0) node[right] {$y$};
	\draw[axis] (0,0,-1) -- (0,0,10) node[left] {$z$};

	\draw[opacity=.5,very thick,fill=Red]
        (2,2,2) -- (2,2,8) -- (\a,\a,8) -- (\a,\a,2) -- cycle;
	\draw[opacity=.5,very thick,fill=Blue]
        (\a,2,2) -- (2,\a,2) -- (2,\a,8) -- (\a,2,8) -- cycle;
	\draw[opacity=.5,very thick,fill=Yellow]
        (2,2,5) -- (2,\a,5) -- (\a,\a,5) -- (\a,2,5) -- cycle;
    \draw[very thick, dashed] (4.12132034356,4.12132034356,2) --
        (4.12132034356,4.12132034356,8);
    \draw[fill] (4.12132034356,4.12132034356,5) circle (3pt);
\end{tikzpicture}
\end{minipage}
\hfill
\begin{minipage}{0.48\textwidth}
\tdplotsetmaincoords{80}{100}
\begin{tikzpicture}[
        tdplot_main_coords,
        axis/.style={->,thick},
        scale=0.40
    ]
    \pgfmathsetmacro{\a}{1.40192378865}
    \pgfmathsetmacro{\b}{6.59807621135}

    \foreach \x in {0,1,...,9}
        \foreach \y in {0,1,...,9} {
			\draw[gray,very thin]
                (\x,0,0) -- (\x,9.5,0);
			\draw[gray,very thin] (0,\y,0) -- (9.5,\y,0);
		}
    \foreach \y in {0,1,...,9}
        \foreach \z in {0,1,...,9} {
			\draw[gray,very thin]
                (0,\y,0) -- (0,\y,9.5);
			\draw[gray,very thin] (0,0,\z) -- (0,9.5,\z);
		}
    \foreach \z in {0,1,...,9}
        \foreach \x in {0,1,...,9} {
			\draw[gray,very thin]
                (0,0,\z) -- (9.5,0,\z);
			\draw[gray,very thin] (\x,0,0) -- (\x,0,9.5);
		}

	\draw[axis] (-2,0,0) -- (10,0,0) node[left] {$x$};
	\draw[axis] (0,-1,0) -- (0,10,0) node[right] {$y$};
	\draw[axis] (0,0,-1) -- (0,0,10) node[left] {$z$};

	\draw[opacity=.5,very thick,fill=Red]
        (2,2,4) -- (8,2,4) -- (8,8,4) -- (2,8,4) -- cycle;
	\draw[opacity=.5,very thick,fill=Blue]
        (2,3.5,\a) -- (8,3.5,\a) -- (8,6.5,\b) -- (2,6.5,\b) -- cycle;
	\draw[opacity=.5,very thick,fill=Yellow]
        (2,3.5,\b) -- (8,3.5,\b) -- (8,6.5,\a) -- (2,6.5,\a) -- cycle;
    \draw[very thick] (2,5,4) -- (8,5,4);
\end{tikzpicture}
\end{minipage}

Wow, not going to lie, it took some time to draw the diagrams above.
Anyhow, with that in mind, let's discuss other fundamental terms
and concepts before we move on to the next section of the note.

\begin{definition}
A linear system is \textbf{consistent} if it retains at least one
solution, whereas a system is \textbf{inconsistent} if there is no
solution.
\end{definition}

We could also define a system based on the column matrix $b$.

\begin{definition}
A linear system is \textbf{homogeneous} if the column matrix $b$
in the linear system $Ax = b$ is a zero matrix. If the matrix is not
a zero matrix, the system is \textbf{nonhomogeneous}, or
\textbf{inhomogeneous}.
\end{definition}

Now from the definitions, we can see that all homogeneous systems
are consistent. As you have guessed, we have a special name for
a special solution in such cases.

\begin{definition}
\textbf{Trivial solution} is a solution where all the entries for
$x$ in $Ax = b$ are zero.
\end{definition}

With the fundamental terms in mind, let's take a look at how matrices
can be used to solve complex linear system with Gaussian Elimination.

\section{Gaussian Elimination}

When you solve a linear system, you probably multiplied constants
and add or subtract equations to cancel variables. That is exactly
what we will be doing in Gaussian Elimination, but with matrices.
First, define how we solve a linear system in a more rigorous way.

\begin{definition}
For a system of linear equations, performing following operations
will result in a linear system with identical solution.
\begin{enumerate}[noitemsep]
\item
Change the order of linear equations in which they are listed
\item
Multiply a nonzero constant, or nonzero scalars, to an equation
\item
Add equations from the system to form a new equation
\end{enumerate}
\end{definition}

Now before we connect this idea with matrices, let's define a matrix
that we will use with the property of a linear system above.

\begin{definition}
An \textbf{augmented matrix} for a linear system $Ax = b$ is a
partitioned matrix in the form $[A \mid b]$.
\end{definition}

Consider the following system for an example.
\begin{align*}
    x + 2y + 3z &= 4 \\
    2x + 3y + 4z &= 5 \\
    3x + 4y + 5z &= 6
\end{align*}
In the linear system above, the augmented matrix will be the following
matrix.
\[
\begin{bNiceArray}{ccc|c}
    1 & 2 & 3 & 4 \\
    2 & 3 & 4 & 5 \\
    3 & 4 & 5 & 6
\end{bNiceArray}
\]
Now that we have augmented matrix in mind, let's take a look at
how the three properties that we discussed applies with matrices.

\begin{definition}
For an augmented matrix for a linear system, performing the following
operations will result in a linear system with identical solution as
the initial system.
\begin{enumerate}[noitemsep]
\item[$R_1$.]
Rearrangement of the rows in the augmented matrix
\item[$R_2$.]
Multiply a nonzero scalar to a row in the augmented matrix
\item[$R_3$.]
Add a row multiplied by a scalar with another to replace a row with
new row
\end{enumerate}
Such operations are referred to as \textbf{elementary row operations}.
\end{definition}

If you compare the operations one by one, you can notice that the
two set of properties refer to the same property. Now with the
set of property of augmented matrices in mind, here are the steps
for the Gaussian Elimination algorithm.

\begin{theorem}{Gaussian Elimination}
Execute the following steps in order for Gaussian Elimination.
\begin{enumerate}[noitemsep]
\item
Construct the augmented matrix.
\item
Utilize the elementary row operations to transform the augmented
matrix into REF.
\item
Convert the augmented matrix in REF into equations of a transformed
linear system, referred to as the \textbf{derived set}.
\item
Back substitute to solve the system.
\end{enumerate}
\end{theorem}

Let's take a look at few examples.

\begin{problem}
Find solution(s) to the following system of linear equations if any.
\begin{align*}
    x + 2y + z &= 1 \\
    3x - y + z &= 6 \\
    x + y - z &= -2
\end{align*}
\end{problem}
\begin{solution}
First, the augmented matrix can be found from the system. Consider
the matrix below.
\[
\begin{bNiceArray}{ccc|c}
    1 & 2  & 1  & 1 \\
    3 & -1 & 1  & 6 \\
    1 & 1  & -1 & -2
\end{bNiceArray}
\]
Using the elementary row operations, the augmented matrix can be
transformed as the following.
\[
\begin{bNiceArray}{ccc|c}
    1 & 2  & 1  & 1 \\
    3 & -1 & 1  & 6 \\
    1 & 1  & -1 & -2
\end{bNiceArray}
\rightarrow
\begin{bNiceArray}{ccc|c}
    1 & 2  & 1  & 1 \\
    0 & -4 & 4  & 12 \\
    1 & 1  & -1 & -2
\end{bNiceArray}
\rightarrow
\begin{bNiceArray}{ccc|c}
    1 & 2  & 1  & 1 \\
    0 & -4 & 4  & 12 \\
    0 & -1 & -2 & -3
\end{bNiceArray}
\]
\[
\rightarrow
\begin{bNiceArray}{ccc|c}
    1 & 2  & 1  & 1 \\
    0 & -1 & 1  & 3 \\
    0 & 1  & 2  & 3
\end{bNiceArray}
\rightarrow
\begin{bNiceArray}{ccc|c}
    1 & 2  & 1  & 1 \\
    0 & -1 & 1  & 3 \\
    0 & 0  & 3  & 6
\end{bNiceArray}
\rightarrow
\begin{bNiceArray}{ccc|c}
    1 & 2  & 1  & 1 \\
    0 & 1  & -1 & -3 \\
    0 & 0  & 1  & 2
\end{bNiceArray}
\]
Because the augmented matrix is now in REF, the transformed
linear system can be rewritten.
\begin{align*}
    x + 2y + z &= 1 \\
    y - z &= -3 \\
    z &= 2
\end{align*}
Back substituting, it is evident that $z = 2$ and $y = -1$.
Moreover, $x = 1 + 2 - 2 = 1$. Therefore, the triple $(1, -1, 2)$
is the solution to the linear system.
\end{solution}

Try substituting to the original system! You can notice that the
solution is valid. Below is the second example of using
Gaussian Elimination.

\begin{problem}
Find solution(s) to the following system of linear equations if any.
\begin{align*}
    x + 2y + 3z &= 4 \\
    y + 3z &= 5 \\
    x + y &= 1
\end{align*}
\end{problem}
\begin{solution}
First and foremost, the augmented matrix can be found.
\[
\begin{bNiceArray}{ccc|c}
    1 & 2 & 3 & 4 \\
    0 & 1 & 3 & 5 \\
    1 & 1 & 0 & 1
\end{bNiceArray}
\rightarrow
\begin{bNiceArray}{ccc|c}
    1 & 2  & 3  & 4 \\
    0 & 1  & 3  & 5 \\
    0 & -1 & -3 & -3
\end{bNiceArray}
\rightarrow
\begin{bNiceArray}{ccc|c}
    1 & 2 & 3 & 4 \\
    0 & 1 & 3 & 5 \\
    0 & 0 & 0 & 2
\end{bNiceArray}
\]
Writing the transformed augmented matrix in to a system of linear
equations, the following system is obtained.
\begin{align*}
    x + 2y + 3z &= 4 \\
    y + 3z &= 5 \\
    0 &= 2
\end{align*}
Notice that $0 = 2$ can never be true. Therefore, no solution exists.
\end{solution}

We could double check out answer without using matrices. Notice that
$3z - x = 4$ and $x = 3z - 4$. Similarly, $y = 5 - 3z$. Substituting
to the first equation, $3z - 4 + 10 - 6z + 3z = 4$, or $6 = 4$,
which is not true. Therefore, our results align! Below is the third
example of using the algorithm.

\begin{problem}
Find solution(s) to the following system of linear equations if any.
\begin{align*}
    x - y - 2z &= 3 \\
    2x + y - 7z &= 3 \\
    x + y - 4z &= 1
\end{align*}
\end{problem}
\begin{solution}
Consider the following augmented matrix.
\[
\begin{bNiceArray}{ccc|c}
    1 & -1 & -2 & 3 \\
    2 & 1  & -7 & 3 \\
    1 & 1  & -4 & 1
\end{bNiceArray}
\]
Utilizing Gaussian Elimination, the augmented matrix can be
transformed as following.
\begin{align*}
\begin{bNiceArray}{ccc|c}
    1 & -1 & -2 & 3 \\
    2 & 1  & -7 & 3 \\
    1 & 1  & -4 & 1
\end{bNiceArray}
&\rightarrow
\begin{bNiceArray}{ccc|c}
    1 & -1 & -2 & 3 \\
    2 & 1  & -7 & 3 \\
    0 & 2  & -2 & -2
\end{bNiceArray}
\rightarrow
\begin{bNiceArray}{ccc|c}
    1 & -1 & -2 & 3 \\
    0 & 1  & -1 & -1 \\
    2 & 1  & -7 & 3
\end{bNiceArray} \\
&\rightarrow
\begin{bNiceArray}{ccc|c}
    1 & -1 & -2 & 3 \\
    0 & 1  & -1 & -1 \\
    0 & 3  & -3 & -3
\end{bNiceArray}
\rightarrow
\begin{bNiceArray}{ccc|c}
    1 & -1 & -2 & 3 \\
    0 & 1  & -1 & -1 \\
    0 & 0  & 0  & 0
\end{bNiceArray}
\end{align*}
Rewriting the system, the following equations are obtained.
\begin{align*}
    x - y - 2z &= 3 \\
    y - z &= -1
\end{align*}
Notice that $y = z - 1$ and $x = y + 2z + 3 = z - 1 + 2z + 3 =
3z + 2$. Because there exists infinitely many $z$, the system
has infinitely many solutions.
\end{solution}

When writing such solutions, it could be written as the following.
\[
x =
\begin{bmatrix}
    3z + 2 \\
    z - 1 \\
    z
\end{bmatrix}
\]
For the next section, let's discuss the inverse of a matrix and its
application in linear systems.

\section{The Multiplicative Inverse}

As we discuss matrix multiplication, you may have wondered about
matrix division. I mean it would be so nice if we could just
divide coefficient matrix to find our variables. Although there
is no ``division'' operation per se, but we could achieve similar
operation with multiplicative inverse of a matrix, or simply
inverse matrix. In other words, some solutions to linear systems
can be written as $x = A^{-1}b$.

\begin{definition}
The \textbf{inverse of a matrix} is a matrix that produces identity
matrix when multiplied with the original matrix. The inverse of
$A$ is denoted as $A^{-1}$.
\end{definition}

We could see that matrices $A$ and $B$ are the inverse of each other
if the following equation is satisfied.
\[
AB = BA = I
\]
Below is an example where $B = A^{-1}$.
\[
A =
\begin{bmatrix}
    1 & 2 \\
    3 & 5
\end{bmatrix},
\quad
B =
\begin{bmatrix}
    -5 & 2 \\
    3  & -1
\end{bmatrix}
\]
One thing to note from the definition is that $A$ and $B$ must be
square matrices. Say we have $A_{4 \times 2}$ and $B_{2 \times 4}$.
Even if both $AB$ and $BA$ returns an identity matrix, the order will
be different. Therefore, for a matrix to have the inverse, it must
be a square matrix.

\begin{definition}
A matrix is \textbf{invertible}, or \textbf{nonsingular}, if it
retains an inverse matrix. A matrix is \textbf{singular} if it
does not have an inverse matrix.
\end{definition}

Now from the definition above, we could establish a theorem.

\begin{theorem}{Diagonal Elements and Invertibility}
If all diagonal elements of a matrix are nonzero constants,
then the matrix is invertible.
\end{theorem}
\begin{proof}
Consider the matrix $A_{m \times m}$ defined below.
\[
A =
\begin{bmatrix}
    a_{11} & a_{12} & a_{13} & \cdots & a_{1m} \\
    a_{21} & a_{22} & a_{23} & \cdots & a_{2m} \\
    a_{31} & a_{32} & a_{33} & \cdots & a_{3m} \\
    \vdots & \vdots & \vdots & \ddots & \vdots \\
    a_{m1} & a_{m2} & a_{m3} & \cdots & a_{mm}
\end{bmatrix}
\]
It suffices to show there there exists the inverse matrix
for all values of the elements. Consider the inverse below.
\[
A^{-1} =
\begin{bmatrix}
    \frac{1}{a_{11}} & 0 & 0 & \cdots & 0 \\
    0 & \frac{1}{a_{22}} & 0 & \cdots & 0 \\
    0 & 0 & \frac{1}{a_{33}} & \cdots & 0 \\
    \vdots & \vdots & \vdots & \ddots & \vdots \\
    0 & 0 & 0 & \cdots & \frac{1}{a_{mm}}
\end{bmatrix}
\]
Because $a_{ii}$ for integer $i \in [1, m]$ are nonzero numbers,
the all elements in the matrix $A^{-1}$ are defined, and the
inverse matrix exists.
\end{proof}

Now, it is important to note that not all square matrices have the
inverse. For an example, consider a matrix with zero row. The same
row in the product matrices will become a zero row. By definition,
it is not possible to make an identity matrix with a zero row.

Now that we discussed the existence, we could also investigate if
the inverse matrix is unique just as a nonzero integer has a unique
reciprocal.

\begin{theorem}{Uniqueness of Inverse Matrix}
If there exists an inverse matrix, the inverse is unique.
\end{theorem}
\begin{proof}
For the sake of contradiction, let $B$ and $C$ be different
inverse of $A$. Consider the following equation.
\[
B = BI = B(AC) = (BA)C = IC = C
\]
Because the fact that $B = C$ contradicts with the initial assumption
that $B \neq C$, all inverse matrices are unique.
\end{proof}

Before we discuss how we actually find an inverse, let's discuss
four important properties of inverse matrices.

\begin{theorem}{Properties of the Inverse Matrix}
For invertible square matrices $A$ and $B$ and a nonzero scalar
$\lambda$, the following equations hold.
\begin{enumerate}
\item $(A^{-1})^{-1} = A$
\item $(A^\intercal)^{-1} = (A^{-1})^\intercal$
\item $(AB)^{-1} = B^{-1}A^{-1}$
\item $(\lambda A)^{-1} = (\frac{1}{\lambda})A^{-1}$
\end{enumerate}
\end{theorem}

There's lot to prove and let's do them one by one starting with
the first one!

\begin{proof}
Let matrix $B = A^{-1}$. Substituting, it suffices to show that
$B^{-1} = A$. By definition of inverse matrix, the equation
$AB = BA = 1$ holds. In other words, $A$ is the inverse of $B$
and the property holds.
\end{proof}

Below is the proof for the second property.

\begin{proof}
Notice that $(AB)^\intercal = B^\intercal A^\intercal$ from
Theorem \ref{theo:Properties of Transpose of a Matrix}. Let matrix
$B = A^{-1}$ and consider the following equation.
\[
I
= (AB)^\intercal
= B^\intercal A^\intercal
= \left( A^{-1} \right)^\intercal A^\intercal
\]
In other words, $(A^{-1})^\intercal$ is the inverse of $A^\intercal$,
or $(A^\intercal)^{-1}$, proving the property.
\end{proof}

Continuing, here is the proof for the third property.

\begin{proof}
Consider the following equation that uses the associative property
of matrix multiplication.
\begin{align*}
    (B^{-1}A^{-1})(AB)
    &= B^{-1}(A^{-1}A)B
    = B^{-1}(I)B
    = B^{-1}(IB) \\
    &= B^{-1}B
    = I
\end{align*}
Because $B^{-1}A^{-1}$ is the inverse of $AB$, the property is true.
\end{proof}

Finally, here is the proof for the last property.

\begin{proof}
Let $B = \lambda A$. Consider the following equation.
\[
\frac{1}{\lambda} A^{-1} B
= \frac{1}{\lambda} A^{-1} \lambda A
= A^{-1} A
= I
\]
In other words, $\frac{1}{\lambda} A^{-1}$ is the inverse of
$\lambda A$, proving the fourth property.
\end{proof}

Now that we took a look at different properties of inverse matrices,
let's discuss how to actually find one with tools to consider.

\begin{definition}
For a matrix $A$ that is not necessarily a square matrix, an
\textbf{elementary matrix}, denoted as $E$, is a square matrix that
performs elementary row operation by multiplying with $A$.
\end{definition}

Let's take a look at an example. Consider matrix $A$ defined below,
and we want to perform elementary row operations of changing the
order of rows $1$ and $3$ and add two times the fourth row to the
second respectively.
\[
A =
\begin{bmatrix}
    1 & 2 & 3 & 4 \\
    1 & 0 & 2 & 0 \\
    2 & 4 & 6 & 8 \\
    1 & 1 & 1 & 1
\end{bmatrix}
\]
The elementary row operations above can be complicated with words,
but we can represent that with two elementary matrices respectively.
Letting $E_1$ and $E_2$ be elementary matrices for each operation,
we can represent the elementary row operations as the following.
\[
E_1 =
\begin{bmatrix}
    0 & 0 & 1 & 0 \\
    0 & 1 & 0 & 0 \\
    1 & 0 & 0 & 0 \\
    0 & 0 & 0 & 1
\end{bmatrix},
\quad
E_2 =
\begin{bmatrix}
    1 & 0 & 0 & 0 \\
    0 & 1 & 0 & 2 \\
    0 & 0 & 1 & 0 \\
    0 & 0 & 0 & 1
\end{bmatrix}
\]
Let's multiply to see if we are right.
\[
E_1 A =
\begin{bmatrix}
    \textcolor{Red}{2} &
    \textcolor{Red}{4} &
    \textcolor{Red}{6} &
    \textcolor{Red}{8} \\
    1 & 0 & 2 & 0 \\
    \textcolor{Red}{1} &
    \textcolor{Red}{2} &
    \textcolor{Red}{3} &
    \textcolor{Red}{4} \\
    1 & 1 & 1 & 1
\end{bmatrix},
\quad
E_2 A =
\begin{bmatrix}
    1 & 2 & 3 & 4 \\
    \textcolor{Red}{3} &
    \textcolor{Red}{2} &
    \textcolor{Red}{4} &
    \textcolor{Red}{2} \\
    2 & 4 & 6 & 8 \\
    1 & 1 & 1 & 1
\end{bmatrix}
\]
Now we can notice three observations and tips for constructing
elementary matrices. For each elementary row operations, we could
do the following.

\begin{definition}
\textbf{Methods to construct the elementary matrix for each
elementary row operation:}
\begin{itemize}[noitemsep]
\item
To interchange row $a$ and $b$, write the corresponding identity
matrix and interchange row $a$ and $b$ from the identity matrix.
\item
To multiply a nonzero scalar $k$ to $a^\text{th}$ row, start from the
corresponding identity matrix and change the $1$ from $a^\text{th}$
row of the identity matrix to $k$.
\item
To add row $a$ multiplied by the nonzero scalar $k$ to row $b$,
start from the corresponding identity matrix and find the
$b^\text{th}$ row in the matrix. Then, change the zero from the
$a^\text{th}$ column of the row to $k$.
\end{itemize}
\end{definition}

From this, we can notice the following theorem.

\begin{theorem}{Invertibility of Elementary Matrices}
All elementary matrices are invertible.
\end{theorem}
\begin{proof}
The proof is rather rather simple. Because all elementary row
operations can be inverted with the inverse operation,
all elementary matrices are invertible.
\end{proof}

Building onto the theorem, we could prove the uniqueness.

\begin{theorem}{Uniqueness of Elementary Matrices}
For each elementary row operations, the corresponding elementary
matrix is unique.
\end{theorem}
\begin{proof}
For the sake of contradiction, assume there exist elementary
matrices $E$ and $F \neq E$. By definition, $EA = FA$ and
$A = EE^{-1}A = E^{-1}FA$. By Theorem \ref{theo:Uniqueness of Identity
Matrices}, $E^{-1}F = I$ and $E = F$ is obtained by Theorem
\ref{theo:Uniqueness of Inverse Matrix}. By contradiction, it is
shown that elementary matrices are unique.
\end{proof}

There are more interesting lemmas and theorems, but let's discuss
them in latter notes. However, let's discuss few theorems for each
elementary row operation before moving on to finding an inverse
matrix.

\begin{theorem}{Three Properties of Elementary Matrices}
For each elementary matrix for corresponding elementary row
operations, the following properties hold.
\begin{enumerate}[noitemsep]
\item
For an elementary matrix that changes the order of two rows,
the inverse of such matrices is itself.
\item
For an elementary matrix that scales a row with nonzero constant
$k$, the inverse can be obtained by replacing $k$ in the elementary
matrix with $\frac{1}{k}$.
\item
For a elementary matrix that adds the scaled row into the other,
its inverse can be obtained by replacing the nonzero scalar $k$
with $-k$.
\end{enumerate}
\end{theorem}

As always, let's start by proving the first property.

\begin{proof}
The inverse of an elementary matrix can be considered as undoing
the elementary row operation. To undo the interchanging of two rows,
the two rows can be interchanged again. In other words, because
the elementary row operation $EE = I$ by Theorem \ref{theo:Uniqueness
of Identity Matrices}, $E = E^{-1}$.
\end{proof}

Continuing, here is the proof for the second property.

\begin{proof}
Let elementary matrices $E$ and $F$ be matrices that scales a
row by a nonzero scalar $k$ and $\frac{1}{k}$ respectively.
By definition, $F(EI) = I$ and $FE = I$. Therefore, $F = E^{-1}$
and the property is shown.
\end{proof}

Lastly, here is the proof for the final property.

\begin{proof}
Similar to the proof for the second property, let elementary matrices
$E$ and $F$ represents the elementary row operations that adds a
row scaled by nonzero constant $k$ to another and the operation
that adds a row scaled by $-k$ to another respectively. By definition,
$F(EI) = I$ and $F = E^{-1}$ as demonstrated above.
\end{proof}

Now before we actually find an inverse, we need to know which
square matrices have the inverse matrix.

\begin{definition}
A square matrix $A$ is invertible if and only if it can be represented
as identity matrix after sequences of elementary row operations, i.e.
there exists elementary matrices $E_i$ such that the following
equation is true.
\[
E_k E_{k-1} \cdots E_1 A = I
\]
\end{definition}

This is a very important result that we will prove later in a note,
so don't worry. Finally, here are the steps to find an inverse if
there exists one.

\begin{definition}
\textbf{Steps to Find the Inverse Matrix of
$\mathbf{A_{m \times m}}$:}
\begin{enumerate}[noitemsep]
\item
Write an augmented matrix $[A \mid I]$ for $I_{m}$.
\item
Utilize elementary row operations on the augmented matrix for
the block $A$ to be in REF.
\item
If the main diagonal of the left partition contains zero,
$A$ is singular. If not, the next step can be continued.
\item
Use elementary row operations again to transform the block $A$
in REF to an identity matrix. The right partition block is
the inverse of $A$.
\end{enumerate}
\end{definition}

Let's take a look at a few examples.

\begin{problem}
Find the inverse of the following matrix $A$.
\[
A =
\begin{bmatrix}
    1 & 2 & 3 \\
    4 & 5 & 6 \\
    7 & 8 & 9
\end{bmatrix}
\]
\end{problem}
\begin{solution}
First, the augmented matrix can be constructed.
\[
\begin{bNiceArray}{ccc|ccc}
    1 & 2 & 3 & 1 & 0 & 0 \\
    4 & 5 & 6 & 0 & 1 & 0 \\
    7 & 8 & 9 & 0 & 0 & 1
\end{bNiceArray}
\]
Consider the following elementary row operations.
\begin{align*}
\begin{bNiceArray}{ccc|ccc}
    1 & 2 & 3 & 1 & 0 & 0 \\
    4 & 5 & 6 & 0 & 1 & 0 \\
    7 & 8 & 9 & 0 & 0 & 1
\end{bNiceArray}
&\rightarrow
\begin{bNiceArray}{ccc|ccc}
    1 & 2 & 3 & 1 & 0 & 0 \\
    8 & 10 & 12 & 0 & 2 & 0 \\
    7 & 8 & 9 & 0 & 0 & 1
\end{bNiceArray} \\
&\rightarrow
\begin{bNiceArray}{ccc|ccc}
    1 & 2 & 3 & 1 & 0 & 0 \\
    1 & 2 & 3 & 0 & 2 & -1 \\
    7 & 8 & 9 & 0 & 0 & 1
\end{bNiceArray} \\
&\rightarrow
\begin{bNiceArray}{ccc|ccc}
    1 & 2 & 3 & 1 & 0 & 0 \\
    0 & 0 & 0 & -1 & 2 & -1 \\
    7 & 8 & 9 & 0 & 0 & 1
\end{bNiceArray}
\end{align*}
It is evident that the main diagonal will contain zero after
the transformation. Therefore, the matrix is not invertible.
\end{solution}

\begin{problem}
Find the inverse of the following matrix $A$.
\[
A =
\begin{bmatrix}
    1 & 0 & 0 \\
    2 & 1 & 1 \\
    0 & 1 & 2
\end{bmatrix}
\]
\end{problem}
\begin{solution}
First and foremost, the augmented matrix can be constructed.
\begin{align*}
\begin{bNiceArray}{ccc|ccc}
    1 & 0 & 0 & 1 & 0 & 0 \\
    2 & 1 & 1 & 0 & 1 & 0 \\
    0 & 1 & 2 & 0 & 0 & 1
\end{bNiceArray}
&\rightarrow
\begin{bNiceArray}{ccc|ccc}
    1 & 0 & 0 & 1 & 0 & 0 \\
    0 & 1 & 2 & 0 & 0 & 1 \\
    2 & 1 & 1 & 0 & 1 & 0
\end{bNiceArray} \\
&\rightarrow
\begin{bNiceArray}{ccc|ccc}
    1 & 0 & 0 & 1 & 0 & 0 \\
    0 & 1 & 2 & 0 & 0 & 1 \\
    0 & 0 & -1 & -2 & 1 & -1
\end{bNiceArray} \\
&\rightarrow
\begin{bNiceArray}{ccc|ccc}
    1 & 0 & 0 & 1 & 0 & 0 \\
    0 & 1 & 2 & 0 & 0 & 1 \\
    0 & 0 & 1 & 2 & -1 & 1
\end{bNiceArray}
\end{align*}
Continuing, the left partition can be transformed into an
identity matrix with elementary row operations.
\[
\begin{bNiceArray}{ccc|ccc}
    1 & 0 & 0 & 1 & 0 & 0 \\
    0 & 1 & 2 & 0 & 0 & 1 \\
    0 & 0 & 1 & 2 & -1 & 1
\end{bNiceArray}
\rightarrow
\begin{bNiceArray}{ccc|ccc}
    1 & 0 & 0 & 1 & 0 & 0 \\
    0 & 1 & 0 & -4 & 2 & -1 \\
    0 & 0 & 1 & 2 & -1 & 1
\end{bNiceArray}
\]
Therefore, the inverse of $A$ is the following.
\[
A^{-1} =
\begin{bmatrix}
    1  & 0  & 0 \\
    -4 & 2  & -1 \\
    2  & -1 & 1
\end{bmatrix}
\]
\end{solution}

We could actually multiply to see if we get $I_3$.
\[
\begin{bmatrix}
    1 & 0 & 0 \\
    2 & 1 & 1 \\
    0 & 1 & 2
\end{bmatrix}
\begin{bmatrix}
    1  & 0  & 0 \\
    -4 & 2  & -1 \\
    2  & -1 & 1
\end{bmatrix}
=
\begin{bmatrix}
    1 & 0 & 0 \\
    0 & 1 & 0 \\
    0 & 0 & 1
\end{bmatrix}
\]
Indeed it is the inverse! In the next section of the note,
we will discuss LU decomposition before completing our note on
linear systems.

\section{LU Decomposition}

The last popular method to solve a linear system that we will
discuss is LU decomposition. As always, let's start by discussing
fundamental definitions.

\begin{definition}
\textbf{LU decomposition} is the process of factorizing an
invertible matrix into the product of lower and upper triangular
matrices where the diagonal elements of the lower triangular matrix
is $1$.
\end{definition}

Before we look at an important theorem and how it can be used to
solve a linear system, here is an example of LU decomposition
for matrix $A$ defined as the following.
\[
A =
\begin{bmatrix}
    1 & 2  & 3 \\
    2 & 6  & 9 \\
    3 & 10 & 18
\end{bmatrix}
\]
From here, notice that matrix $A$ can transform to upper triangular
matrix with following elementary row operations of $R_3$.
\[
\begin{bmatrix}
    1 & 2  & 3 \\
    2 & 6  & 9 \\
    3 & 10 & 18
\end{bmatrix}
\rightarrow
\begin{bmatrix}
    1 & 2  & 3 \\
    0 & 2  & 3 \\
    3 & 10 & 18
\end{bmatrix}
\rightarrow
\begin{bmatrix}
    1 & 2 & 3 \\
    0 & 2 & 3 \\
    0 & 4 & 9
\end{bmatrix}
\rightarrow
\begin{bmatrix}
    1 & 2 & 3 \\
    0 & 2 & 3 \\
    0 & 0 & 3
\end{bmatrix}
\]
Notice that the operations above can be represented as the
matrix multiplication with $R_3$ that are lower triangular matrices.
\[
\begin{bmatrix}
    1 & 0  & 0 \\
    0 & 1  & 0 \\
    0 & -2 & 1
\end{bmatrix}
\begin{bmatrix}
    1  & 0 & 0 \\
    0  & 1 & 0 \\
    -3 & 0 & 1
\end{bmatrix}
\begin{bmatrix}
    1  & 0 & 0 \\
    -2 & 1 & 0 \\
    0  & 0 & 1
\end{bmatrix}
A
=
\begin{bmatrix}
    1 & 2 & 3 \\
    0 & 2 & 3 \\
    0 & 0 & 3
\end{bmatrix}
\]
Therefore, the following equation is obtained.
\begin{align*}
    A
    &=
    \left(
    \begin{bmatrix}
        1 & 0 & 0 \\
        2 & 1 & 0 \\
        0 & 0 & 1
    \end{bmatrix}
    \begin{bmatrix}
        1 & 0 & 0 \\
        0 & 1 & 0 \\
        3 & 0 & 1
    \end{bmatrix}
    \begin{bmatrix}
        1 & 0 & 0 \\
        0 & 1 & 0 \\
        0 & 2 & 1
    \end{bmatrix}
    \right)
    \begin{bmatrix}
        1 & 2 & 3 \\
        0 & 2 & 3 \\
        0 & 0 & 3
    \end{bmatrix} \\
    &=
    \begin{bmatrix}
        1 & 0 & 0 \\
        2 & 1 & 0 \\
        3 & 2 & 1
    \end{bmatrix}
    \begin{bmatrix}
        1 & 2 & 3 \\
        0 & 2 & 3 \\
        0 & 0 & 3
    \end{bmatrix}
\end{align*}
$A$ is now represented as the product of lower and upper triangular
matrices! Now that we know how LU decomposition work, let's take
a look at an important lemma on triangular matrices to prove a
theorem on LU decomposition.

\begin{lemma}{Inverse of Triangular Matrices}
The inverse of upper and lower triangular matrices are upper and
lower triangular matrices respectively.
\end{lemma}
\begin{proof}
First and foremost, consider an invertible upper triangular matrix
$A_{m \times m} = [a_{ij}]$ and its inverse $B = [b_{ij}]$.
By definition, its product $C = [c_{ij}]$ can be represented as
$C = [ \sum_{k = 1}^m a_{ij} b_{ji} ]$. Because $c_{ij} = 0$ for
$i \neq j$ and $c_{ij} = 1$ for $i = j$, the values of $b_{ji}$
can be found.

By definition, $a_{ij} = 0$ for $i > j$. Therefore, $c_{ij} = 0$
for $i > j$. Similarly, if $b_{ji} = 0$ for $j > i$, then $c_{ij} = 0$
for $i \neq j$. Continuing, if $b_{ij} = \frac{1}{a_{ij}}$ for
$i = j$, $B$ can be constructed, which is an upper triangular
matrix. Because the inverse matrix is unique, the inverse of an
upper triangular matrix is upper triangular matrix if there exists
one.

Similarly, the fact that inverse of an lower triangular matrix is
lower triangular matrix can be proven. Let $A'_{m \times m} =
[a'_{ij}]$ be a lower triangular matrix and $B' = [b'_{ij}]$ be
its inverse. By definition, their product $C' = [c'_{ij}] =
[ \sum_{k = 1}^m a'_{ij} b'_{ji} ]$ is the identity matrix $I_m$.
Thus, it suffices to show that $c'_{ij} = 0$ for $i \neq j$ and
$c'_{ij} = 1$ for $i = j$.

Note that by definition, $a'_{ij} = 0$ for $i < j$. Therefore,
$c'_{ij} = 0$ for $i < j$. Similarly, if $b'_{ji} = 0$ for $i > j$,
and $b'_{ij} = \frac{1}{a'_{ij}}$ for $i = j$, then the valid
inverse matrix $B$ is constructed. Because $B$ is a lower triangular
matrix and the inverse matrix is unique, the inverse of a lower
triangular matrix is a lower triangular matrix.
\end{proof}

With the lemma in mind, let's prove the following theorem.

\begin{theorem}{Condition for LU Decomposition}
An invertible matrix has a LU decomposition if and only if
multiplying sequence of lower triangular elementary matrices for
elementary row operations can transform the matrix into an upper
triangular matrix.
\end{theorem}
\begin{proof}
First, the statement that a nonsingular matrix has a LU decomposition
if multiplying sequence of lower triangular elementary matrices for
elementary row operations can transform the matrix into an upper
triangular matrix can be proven. Let $A$ be an invertible matrix
that can be represented as the following.
\[
E_{n,n-1} \cdots E_{31} E_{21} A = U
\]
for lower matrices $E$ and an upper matrix $U$. Notice that $A$ can
be written as the following.
\[
A = \left( E^{-1}_{21} E^{-1}_{31} \cdots E^{-1}_{n,n-1} \right) U
\]
From Lemma \ref{lem:Inverse of Triangular Matrices}, the inverse of
the elementary matrices are lower triangular matrices. Moreover,
by Theorem \ref{theo:Property of Triangular Matrix Multiplication},
the product of the lower triangular matrices is also a lower
triangular matrix. Therefore, for a lower triangular matrix $L$,
$A$ can be written as the following.
\[
A = LU
\]
Therefore, the first half of the statement holds.

For the second half of the statement, consider the following equation.
\[
I = AA^{-1} = LUA^{-1} = L \left( UA^{-1} \right)
\]
Therefore, $L$ is nonsingular. From the proof for Lemma
\ref{lem:Inverse of Triangular Matrices}, it is evident that the diagonal
elements of an invertible lower triangular matrix is nonzero as
the diagonal elements in the inverse cannot be defined if a zero is
present in the diagonal elements of the lower triangular matrix.

Let $D$ be the diagonal matrix formed with the diagonal elements of
$L$ and define $L'$ such that the equation $L'D = L$ is satisfied.
Moreover, let $U' = DU$. Thus, the following equation is obtained.
\[
A = LU = L'DU = L'U'
\]
Recall that $L$ is invertible. Because $L^{-1}$ exists and
$(L'D)^{-1} = D^{-1} L'^{-1}$ exist, $L'$ is also invertible.
In other words, $L'^{-1} A = U'$. Therefore, it suffices to show
that $L'^{-1}$ is a product of lower triangular elementary matrices.
By definition of $D$, the diagonal elements of $L'$ are $1$. Thus,
it can be interpreted as a sequence of elementary row operations.
In other words, it can be represented as a product of lower triangular
elementary matrices, proving the latter half of the theorem.
\end{proof}

Now that we discussed what LU decomposition is, let's see how
we can use it to solve a linear system.

\subsection{Application in Linear Systems}

When we apply LU decomposition in solving linear systems, two
relatively easy systems must be solved. Consider the system
$Ax = b$ where $A = LU$. First, $y$ must be found from the
following system with forward substitution.
\[
Ly = b
\]
Then, we can solve for $x$ with back substitution with the following
system.
\[
Ux = y
\]
The good news is that both system is relatively easy to solve.
The reason is simple. Consider the following equation.
\[
Ax = LUx = b
\]
Let $y = Ux$, then $Ly = b$. If we find $Ly = b$, then we can find
$Ux = y$ as we know $y$. This way, we can find $x$ while exploiting
the triangular form. Let's take a look at an example.

\begin{problem}
Solve the following system of linear equations.
\begin{align*}
    x - 2y + z &= 2 \\
    2x + y - z &= 1 \\
    3x - y + 2z &= 15
\end{align*}
\end{problem}
\begin{solution}
First and foremost, the linear system can be represented as the
matrix multiplication below.
\[
\begin{bmatrix}
    1 & -2 & 1 \\
    2 & 1  & -1 \\
    3 & -1 & 2
\end{bmatrix}
\begin{bmatrix}
    x \\
    y \\
    z
\end{bmatrix}
=
\begin{bmatrix}
    2 \\
    1 \\
    15
\end{bmatrix}
\]
Continuing, the coefficient matrix $A$ can be transformed as the
following.
\begin{align*}
    \begin{bmatrix}
        1 & -2 & 1 \\
        2 & 1  & -1 \\
        3 & -1 & 2
    \end{bmatrix}
    &\rightarrow
    \begin{bmatrix}
        1 & -2 & 1 \\
        0 & 5  & -3 \\
        3 & -1 & 2
    \end{bmatrix}
    \rightarrow
    \begin{bmatrix}
        1 & -2 & 1 \\
        0 & 5  & -3 \\
        0 & 5  & -1
    \end{bmatrix} \\
    &\rightarrow
    \begin{bmatrix}
        1 & -2 & 1 \\
        0 & 5  & -3 \\
        0 & 0  & 2
    \end{bmatrix}
\end{align*}
Using elementary matrices, the transformation above can be
written as the following.
\[
\begin{bmatrix}
    1 & 0  & 0 \\
    0 & 1  & 0 \\
    0 & -1 & 1
\end{bmatrix}
\begin{bmatrix}
    1  & 0 & 0 \\
    0  & 1 & 0 \\
    -3 & 0 & 1
\end{bmatrix}
\begin{bmatrix}
    1  & 0 & 0 \\
    -2 & 1 & 0 \\
    0  & 0 & 1
\end{bmatrix}
A =
\begin{bmatrix}
    1 & -2 & 1 \\
    0 & 5  & -3 \\
    0 & 0  & 2
\end{bmatrix}
\]
Completing the LU decomposition, $A$ can be represented as the
product of following lower and upper triangular matrices.
\begin{align*}
    A
    &=
    \left(
    \begin{bmatrix}
        1 & 0 & 0 \\
        2 & 1 & 0 \\
        0 & 0 & 1
    \end{bmatrix}
    \begin{bmatrix}
        1 & 0 & 0 \\
        0 & 1 & 0 \\
        3 & 0 & 1
    \end{bmatrix}
    \begin{bmatrix}
        1 & 0 & 0 \\
        0 & 1 & 0 \\
        0 & 1 & 1
    \end{bmatrix}
    \right)
    \begin{bmatrix}
        1 & -2 & 1 \\
        0 & 5  & -3 \\
        0 & 0  & 2
    \end{bmatrix} \\
    &=
    \begin{bmatrix}
        1 & 0 & 0 \\
        2 & 1 & 0 \\
        3 & 1 & 1
    \end{bmatrix}
    \begin{bmatrix}
        1 & -2 & 1 \\
        0 & 5  & -3 \\
        0 & 0  & 2
    \end{bmatrix}
\end{align*}
Continuing, the following system could be solved.
\[
\begin{bmatrix}
    1 & 0 & 0 \\
    2 & 1 & 0 \\
    3 & 1 & 1
\end{bmatrix}
\begin{bmatrix}
    x' \\
    y' \\
    z'
\end{bmatrix}
=
\begin{bmatrix}
    2 \\
    1 \\
    15
\end{bmatrix}
\]
The matrix multiplication above can be written as the following
system of linear equations.
\begin{align*}
    x' &= 2 \\
    2x' + y' &= 1 \\
    3x' + y' + z' &= 15
\end{align*}
Substituting, $y' = 1 - 4 = -3$ and $z' = 15 - 6 + 3 = 12$. Thus,
the final system can be written as the following.
\[
\begin{bmatrix}
    1 & -2 & 1 \\
    0 & 5  & -3 \\
    0 & 0  & 2
\end{bmatrix}
\begin{bmatrix}
    x \\
    y \\
    z
\end{bmatrix}
=
\begin{bmatrix}
    2 \\
    -3 \\
    12
\end{bmatrix}
\]
Writing in linear system, $x$, $y$, and $z$ can be found.
\begin{align*}
    x - 2y + z &= 2 \\
    5y - 3z &= -3 \\
    2z &= 12
\end{align*}
Therefore, $z = 6$, $y = 3$, and $x = 2$.
\end{solution}

With LU decomposition, we completed our discussion for solving
system of linear equations! Below are some practice problems
for the topics that we covered.

\section{Practice Problems}
\begin{enumerate}
\item
Show that for an invertible square matrix $A$, the system $Ax = b$
is always consistent.

\item
Show that the following system cannot have a unique solution for all
$k \in \mathbb{R}$.
\[
\begin{bmatrix}
    1 & 2 \\
    2 & 4
\end{bmatrix}
\begin{bmatrix}
    x \\
    y
\end{bmatrix}
=
\begin{bmatrix}
    1 \\
    k
\end{bmatrix}
\]

\item
Use Gaussian Elimination to find the solution for the following
system.
\begin{align*}
    x + y + z   &= 10 \\
    3x + 4y - z &= 2  \\
    -x - 3y + z &= 4
\end{align*}

\item
Solve the same system above using LU decomposition.

\item
Find the inverse of the following matrix.
\[
\begin{bmatrix}
    1  &  1 & 1 \\
    3  &  4 & -1 \\
    -1 & -3 & 1
\end{bmatrix}
\]
Then, find the values of $x, y, z$ that satisfy the following system.
\[
\begin{bmatrix}
    1  &  1 & 1 \\
    3  &  4 & -1 \\
    -1 & -3 & 1
\end{bmatrix}
\begin{bmatrix}
    x \\
    y \\
    z
\end{bmatrix}
=
\begin{bmatrix}
    10 \\
    2 \\
    4
\end{bmatrix}
\]
\end{enumerate}
\end{document}



